\chapter{Subprocess and PTY embedding}
\label{ch:pty}

\begin{aside}
Sometimes your TUI is a wrapper around another program.
A git-front-end runs git. A debugger driver runs gdb. The
inner program may itself be a TUI; you need a pseudo-
terminal.
\end{aside}

\section{Pipes vs PTYs}

\begin{itemize}[leftmargin=*]
  \item \textbf{Pipe.} One-way byte stream; no terminal
    semantics. The child sees pipe, may behave differently
    (no colour, no line editing).
  \item \textbf{PTY.} Bidirectional; mimics a terminal. The
    child sees a terminal and uses ANSI normally.
\end{itemize}

For TUI children, use a PTY.

\section{Opening a PTY}

\begin{lstlisting}
int master, slave;
char name[64];
openpty(&master, &slave, name, nullptr, nullptr);

pid_t pid = fork();
if (pid == 0) {
    setsid();
    dup2(slave, 0);
    dup2(slave, 1);
    dup2(slave, 2);
    close(master);
    execlp("git", "git", "log", nullptr);
}
close(slave);
\end{lstlisting}

The master fd is for the parent; the slave is the child's
controlling terminal.

\section{Reading from the master}

The child writes ANSI; the parent reads bytes and parses.
\maya{} provides a \code{TerminalView} widget that
maintains a virtual screen and renders it:

\begin{lstlisting}
auto term = terminal_view(master_fd);
\end{lstlisting}

\section{Writing to the child}

Keystrokes go to the master:

\begin{lstlisting}
on_key([&](KeyEvent ev) {
    auto bytes = encode_key(ev);
    write(master_fd, bytes.data(), bytes.size());
});
\end{lstlisting}

\section{Size propagation}

\code{TIOCSWINSZ} on the master notifies the child of size.
Call it whenever the parent terminal resizes.

\section{The parser}

A virtual terminal parser interprets ANSI: maintains a
cell grid, cursor, scrollback. Open-source options:
\code{libvterm}, \code{libtmt}, \code{wezterm}'s
vt-parser. \maya{} can wrap any of these.

\section{Scrollback}

The virtual terminal keeps the last N rows in scrollback.
The widget shows the visible viewport plus scroll
controls.

\section{Embedded vs spawned}

Embedded: the child renders inside your TUI's screen.
Spawned: the child takes over the terminal until it
exits, then your TUI resumes. Use spawned for ``run my
editor for the user to edit a file''.

\section{Recap}

\begin{recap}
\begin{itemize}[leftmargin=*]
  \item PTY for TUI children; pipe for plain programs.
  \item Open with \code{openpty}; fork, exec, dup2.
  \item Virtual terminal parser to render the child.
  \item Propagate resize with \code{TIOCSWINSZ}.
  \item Embedded vs spawned trade-off.
\end{itemize}
\end{recap}

\section*{Workshop}

\begin{enumerate}[leftmargin=*]
  \item Build a terminal-view widget running \code{bash}
    inside your app.
  \item Propagate resize correctly.
  \item Add scrollback with mouse-wheel scroll.
\end{enumerate}
